\chapter{Структурный системный анализ исследуемого объекта}\label{ch:2}

Объектом исследования в данной работе является процесс оценки стоимости ремонта автомобилей в автосервисе. Процесс достаточно обширен и включает в себя множество действий, начиная от обращения клиента и заканчивая принятием решения о начале ремонтных работ (либо отказом). В него вовлечены три основные роли:

\begin{itemize}
    \item Клиент
    \item Сотрудник автосервиса (мастер)
    \item Начальник автосервиса
\end{itemize}

Для погружения в контекст процесса и понимания его деталей, была составлена BPMN-диаграмма, подробно отражающая состояние "as is" в автосервисе. (Рисунок 10)


Проанализировав диаграмму, можно выявить основные проблемные места процесса:
\begin{itemize}
    \item Не все клиенты обращаются в автосервис лично. Многие делают это в режиме онлайн, и этап сбора информации сопряжён с дополнительными трудностями. Иногда может быть недоступен сотрудник автосервиса, к которому обращаются, и клиенту приходится ожидать, пока мастер дойдёт до устройства, обработает обращение и даст обратную связь.
    \item Несогласованный формат отправки изображений мастеру приводит к тому, что иногда информации с изображений недостаточно (например, не понятно, что за деталь), и приходится запрашивать уточнение данных или новое фото. Это приводит к затратам времени на переписку (дополнительные циклы пересъёмки/разъяснений, ожиданий ответа с обеих сторон).
    \item Иногда неопытность сотрудника и отсутствие чёткого понимания расценки работ приводит либо к необходимости согласования с начальником автосервиса (дополнительное время на согласование), либо к тому, что одна и та же поломка может оцениваться по-разному разными мастерами. 
\end{itemize}

\begin{figure}[H]
\centering
\begin{minipage}{0.90\textwidth}
    \centering
    \includegraphics[scale=0.55, angle=90]{BPMN}
\end{minipage}\hspace{-4em}%
\begin{minipage}{0.06\textwidth}
    \centering
    \rotatebox{90}{%
        \parbox{10cm}{\caption{BPMN-диаграмма процесса}}%
    }
\end{minipage}
\end{figure}

По экспертной информации, предоставленной начальником автосервиса, более половины случаев обрабатываются в режиме онлайн, а время обработки одного онлайн-обращения занимает в среднем 30 минут. Это указывает на высокую загруженность мастеров в дни обработки обращений сразу от множества клиентов. Около 20\% обращений обсуждаются с начальником, а с клиентом контакты осуществляются в среднем 2-3 раза, что также говорит о весомых тратах времени на согласование.

Таким образом, в текущем бизнес-процессе имеется несколько существенных недостатков:
\begin{itemize}
    \item Длительные согласования и ожидания ответа
    \item Отвлечение мастеров от работы на приём обращений
    \item Ошибки при оценке стоимости у новых сотрудников
\end{itemize}

По имеющейся статистике, вероятность существенного изменения (более, чем на 5-10\%) предварительной оценки стоимости ремонта - всего около 5\%. Это указывает на устойчивость и воспроизводимость результатов предварительной оценки.